\documentclass[12pt]{article}
\usepackage[utf8]{inputenc}
\usepackage[T1]{fontenc}
\usepackage[margin=1in]{geometry}
\usepackage{setspace}
\usepackage{xcolor}
\usepackage[colorlinks=true, urlcolor=blue!60!black, linkcolor=blue!60!black]{hyperref}
\setlength{\parindent}{0em}
\setlength{\parskip}{0.6em}
\renewcommand{\baselinestretch}{1.0}
\begin{document}
\begin{center}
{\LARGE \textit{Church of the Lukumi Babalu Aye, Inc. v. City of Hialeah}, 508 U.S. 520 (1993)\par}
\vspace{0.5em}
{\large Case Brief --- Two-Page Edition\par}
\end{center}

\section{Memory Jogger}
Hialeah's ordinances, gerrymandered to ban Santeria animal sacrifice while exempting nearly all secular and kosher killing, were neither neutral nor generally applicable under \textit{Smith} and failed strict scrutiny (\href{https://supreme.justia.com/cases/federal/us/508/520/}{508 U.S. 520}).

\section{Detailed Case Facts}
Santeria, an Afro-Caribbean religion brought to South Florida by Cuban exiles, centers on animal sacrifice. When the Church of the Lukumi Babalu Aye announced plans to open a house of worship in Hialeah in 1987, an openly hostile city council passed resolutions and four ordinances banning ritual ``sacrifice'' of animals ``not for the primary purpose of food consumption,'' while exempting licensed slaughterhouses, kosher slaughter, hunting, fishing, and pest control.

\section{Procedural History}
The District Court, though finding the ordinances not religiously neutral, upheld them as serving compelling interests; the Eleventh Circuit affirmed \textit{per curiam}. The Supreme Court reversed.

\section{Judicial Votes}
Reversed, 9--0. Kennedy wrote for the Court, but his legislative-motive section drew only Stevens; Scalia (with Rehnquist) and Souter concurred in part and in the judgment; Blackmun (with O'Connor) concurred in the judgment.

\section{Holding}
A law that targets religious practice, overtly or through a ``religious gerrymander,'' or that pursues its interests only against religious conduct, is not neutral or generally applicable and must survive strict scrutiny. These ordinances could not.

\section{Analysis of Opinions}
Kennedy held that facial neutrality is only the minimum. The interlocking definitions and exemptions left Santeria sacrifice virtually the only conduct banned, and the ordinances were badly underinclusive as to both animal cruelty and public health, so they were neither narrowly tailored nor backed by compelling interests. Scalia refused to probe legislators' motives, focusing on the laws' effects. Souter and Blackmun urged reconsidering \textit{Smith} itself; Blackmun said a targeting law fails strict scrutiny \textit{ipso facto}.

\section{Examples: Future Applications}
\textbf{Same-Side:} a ban on ``ceremonial animal killing'' aimed at a mosque's Eid slaughter, exempting hunting and abattoirs, is unconstitutional. \textbf{Opposite-Side:} a uniform rule that everyone seal organic animal waste is neutral and generally applicable. \textbf{Fence-Sitter:} a neutral anti-discrimination law enforced by commissioners who disparage a baker's faith, as in \textit{Masterpiece Cakeshop}.

\section{Critique}
By not defining how much underinclusiveness defeats general applicability, \textit{Lukumi} left lower courts divided until \textit{Fulton} and \textit{Tandon}, and Part II-A-2's two votes left the role of officials' hostility unsettled. Originalist critics of \textit{Smith} say it protects minorities only from openly hostile majorities, while progressives warn it has become a deregulatory sword.

\section{Key Quotations}
\begin{quote}
``Official action that targets religious conduct for distinctive treatment cannot be shielded by mere compliance with the requirement of facial neutrality. The Free Exercise Clause protects against governmental hostility which is masked as well as overt.''
\end{quote}
\noindent --- Justice Kennedy, Opinion of the Court, 508 U.S. at 534.

\end{document}
